\honest{Anchoring and Adjustment}{Eliezer Yudkowsky}{2007}

I spin a wheel of fortune, it stops at 65, and I ask you what share of the countries in the United Nations are in Africa. Then I ask you to estimate 1 × 2 × 3 × 4 × 5 × 6 × 7 × 8 in five seconds.

Tversky and Kahneman found that people who saw the wheel stop at 65 gave a median estimate of 45\%, and people who saw 10 gave 25\%. \nb{The subjects had been told to reach their estimate ``by moving upward or downward from the given number.''} Students shown the product in ascending order estimated 512, and students shown it in descending order estimated 2,250. The answer is 40,320.

``The current theory'' is that people start from the anchor and adjust until an answer ``sounds plausible,'' and then stop. Numbers further away could also be plausible, so the adjustment usually falls short. In the multiplication task, students presumably multiplied the first few numbers and adjusted upward; both groups fell short, and the ascending group more so because it started lower. \nb{For the wheel, the research the post cites had moved away from this. Strack and Mussweiler, cited in the post for the Einstein study, argued that anchoring is a kind of semantic priming, and Epley and Gilovich (2001) reported that anchors given by an experimenter appear to work by making anchor-consistent information easier to recall. Epley and Gilovich defended adjustment for anchors people produce themselves, as in the multiplication task. A month later, in ``Priming and Contamination,'' Yudkowsky wrote that ``most anchoring is actually due to contamination, not sliding adjustment.''}

Payoffs for accuracy did not reduce the effect. Anchors as absurd as 1215 or 1992, for the year Einstein first visited the United States, worked as well as 1905 or 1939.

When you negotiate a salary or buy a car, watch out for people who exploit this, and notice when you are ``adjusting'' a figure yourself.

Debiasing has ``generally proved not very effective.'' I suggest two remedies of my own: throw away an implausible anchor rather than ``sliding from'' it, and, since subjects told to avoid anchoring still anchor, also think briefly of an anchor that is clearly wrong in the other direction. \nb{The post reports no test of either. A method tested by Mussweiler, Strack and Pfeiffer (2000) did reduce anchoring: listing reasons why the given anchor is inappropriate.}
